\chapter{Оператор циклу \texttt{while}}

Оператори циклів використовують у випадках, коли в програмі необхідно багаторазово виконати певну послідовність команд. У мові Python основним циклом із умовою повторення (циклом із передумовою) є цикл \texttt{while}.

\vspace{0.3cm}

\begin{pydef}
Синтаксис циклу while
while вираз_умови:
    блок_операторів
\end{pydef}

\vspace{0.3cm}

Блок операторів циклу \texttt{while} повторюється до того часу, поки результатом логічного висловлювання умови повтору є значення \texttt{True}. Вираз умови обчислюється перед кожним повтором циклу, і вихід з циклу (тобто перехід до оператора, наступного за оператором \texttt{while}) здійснюється в тому випадку, якщо вираз умови має значення \texttt{False}.

\begin{py}{Приклад 1. Виведення значень лічильника від 0 до 4}
i = 0
while i < 5:
    print("значення лічильника =", i)
    i += 1  # i = i + 1

print("end")
\end{py}

\noindent\textbf{Покроковий аналіз виконання:}

Задається початкове значення змінної $i = 0$.

 Перевіряється умова $i < 5$. Оскільки вона виконується ($\text{\texttt{True}}$), відбувається перехід до операторів всередині циклу (виводимо значення $i$ на екран та збільшуємо на 1, тобто $i = 0 + 1 = 1$). 
 
 Знову повертаємось до перевірки умови на початку $i < 5$. Вона виконується $\rightarrow$ виконуються оператори виводу на екран та збільшення $i$, і так до тих пір, поки значення $i$ не стане дорівнювати 5. Тоді оператори всередині циклу вже не будуть виконуватись, виконання перейде до оператора \texttt{print("end")}.
 \subsection*{Різниця між циклами for та while}
 
 \begin{itemize}
 	\item \textbf{Цикл for} використовується тоді, коли кількість повторень відома наперед.  
 	Він перебирає елементи послідовності (наприклад, \texttt{range()}, список, рядок) і виконує тіло циклу для кожного елемента.
 	
 	\item \textbf{Цикл while} застосовується тоді, коли кількість повторень невідома заздалегідь.  
 	Він виконує тіло циклу доти, доки умова залишається істинною (\texttt{True}).  
 	Якщо умова стає хибною (\texttt{False}), цикл завершується.
 \end{itemize}
 
 

\section{ Керування циклом: \texttt{break}, \texttt{continue}, \texttt{pass} та блок \texttt{else}}

\begin{pydef}
Структура керування циклом while
while вираз_умови_test:
    блок_операторів_1
    if умова_1:
        break     # вийти з циклу, пропустивши else
    if умова_2:
        continue  # перейти на початок циклу до 
                  #виразу_умови_test
    блок_операторів_2
else:
    оператори_else # виконуються, якщо вираз_умови 
                   # став False (без break)
\end{pydef}

\begin{itemize}
	\item \textbf{\texttt{break}} "--- вихід з циклу із пропуском блоку \texttt{else}.
	\item \textbf{\texttt{continue}} "--- перехід на початок циклу до перевірки умови.
	\item \textbf{\texttt{pass}} "--- порожня інструкція. Не виконує жодних дій і використовується у випадках, коли синтаксис мови потребує наявності інструкції, але поки що жодних корисних дій виконати не потрібно:
\begin{minted}[fontsize=\small,frame=single,bgcolor=MintCream]{python}
if True:
    pass  # інструкції будуть написані пізніше

while True:
    pass  # порожній нескінченний цикл
\end{minted}
\end{itemize}

\section{ Типові застосування циклу \texttt{while}}

\begin{py}{Приклад 2. Інтерактивна перевірка умов (перевірка стану)}
x = input("Двері відчинені? ")
while x.lower() != "так":
    print("стукаємо")
    x = input("Двері відчинені? ")
print("Заходимо")
\end{py}

\textit{Примітка:} Блок операторів у \texttt{while} може не виконатись жодного разу, якщо вираз-умова від самого початку має значення \texttt{False}.

\begin{py}{Приклад 3. Табулювання функції y = sin(x) до першого від'ємного значення}
import math

x = 0.0
while math.sin(x) >= 0:
    print(f"sin({x:.1f}) = {math.sin(x):.4f}")
    x += 0.1
\end{py}

\begin{py}{Приклад 4. Запит на введення додатнього числа з валідацією}
x = int(input("Введіть додатнє x = "))
while x <= 0:
    x = int(input("Введіть додатнє x = "))
print("Дякую")
\end{py}
Програма буде виводити на екран повідомлення та перевіряти, яке число ввід користувач. Якщо $x<=0$, то повідомлення з запрошенням вводу буде з'являтися знов та знов. Якщо додатнє, то цикл завершиться, бо умова $x<=0$ перестане виконуватись, та виконаэться рядок коду, що йде одрвзу після циклу while. 

\begin{py}{Приклад 5. Обчислення суми введених чисел до введення 0 (стоп-сигнал)}
s = 0.0
x = float(input())
while x != 0:
    s += x
    x = float(input())
print("Сума =", s)
\end{py}

Ініціалізується змінна $s$, яка буде накопичувати суму введених чисел.

Користувач вводить число з клавіатури, яке перетворюється у тип float (дійсне число).

Запускається цикл, який працює доти, доки введене число не дорівнює нулю.

До змінної s додається введене число x.

Користувач вводить наступне число. Якщо це знову не нуль — цикл повторюється.

Коли користувач вводить 0, цикл завершується, і програма виводить загальну суму всіх введених чисел.

\begin{py}{Приклад 6. Обчислення суми нескінченного ряду (розклад exp(x)) з точністю eps}
# Розклад e^x = 1 + x + x^2/2! + ... + x^n/n! + ...
import math
# Початкові дані та параметри
x = 1.0       # Значення x, для якого обчислюємо e^x
eps = 1e-9    # Точність обчислень (10^-9)
s = 0.0       # Змінна для накопичення суми ряду
term = 1.0    # Перший доданок ряду (при n = 0, x^0 / 0! = 1)
i = 1         # Лічильник ітерацій та порядковий номер 
#наступного додатка

# Цикл виконується, поки модуль поточного доданка більший за 
#задану точність eps
while abs(term) > eps:
    s += term           # Додаємо поточний доданок до  суми
    term = term * x / i # Обчислюємо наступний доданок
    i += 1              # Збільшуємо номер ітерації на 1
print(f"Сума s = {s}")
print(f"Ітерацій: {i - 1}")
\end{py}

\textit{Примітка:} Будь-який цикл \texttt{for} можна записати за допомогою \texttt{while}, але зворотне твердження взагалі кажучи невірно, оскільки \texttt{while} призначений для циклів з невідомою заздалегідь кількістю повторень.

\section{Проблема «зациклювання»}

Іноді при неуважному використанні оператора \texttt{while} виникає так зване \textbf{зациклювання} "--- нескінченне повторення тіла циклу. Це може статися з двох основних причин:

\begin{enumerate}[label=\arabic*.]
	\item Змінні, що є частиною висловлювання умови циклу, не змінюють усередині циклу своїх значень або приймають некоректні значення, у результаті чого вираз умови завжди має значення \texttt{True}.
	\item Некоректним є сам вираз умови циклу (наприклад, порівняння дійсних чисел типу \texttt{float} на строгу рівність).
\end{enumerate}

\begin{py}{Приклад зациклювання 1: Нескінченний цикл при перериванні Ctrl+C}
while True:
    print("Натисніть Ctrl+C для зупинки")
\end{py}
while True: — створює нескінченний цикл. Оскільки умова True (істина) завжди залишається True і ніколи не змінюється на False, цикл буде виконуватися вічно.

print("Натисніть Ctrl+C для зупинки") — тіло циклу, яке безперервно виводить це повідомлення в консоль/термінал.

Зациклювання виникає тому, що у коді немає жодної умови виходу (наприклад, оператора break або зміни логічного виразу), тому програма самостійно зупинитися не може.

Щоб примусово перервати виконання нескінченного циклу в терміналі чи консолі, використовують комбінацію клавіш Ctrl + C. Вона надсилає програмі сигнал переривання (KeyboardInterrupt), що викликає зупинку виконання коду.
\begin{py}{Приклад зациклювання 2: Небезпечне порівняння дійсних чисел (float)}
# Некоректний код 
x = 0.0
while x != 1.0:# умова не перестане виконуватись ніколи
    print(f"x = {x:.1f}")
    x += 0.3

# Правильний варіант:
x = 0.0
while x < 1.0:
    print(f"x = {x:.1f}")
    x += 0.3
\end{py}

\section{Складніші приклади та алгоритми}

\begin{py}{Приклад 7. Середнє арифметичне простих чисел до введення 0}
x = int(input())  # Зчитуємо перше ціле число з клавіатури
k = 0             # Лічильник знайдених простих чисел
S = 0             # Змінна для збереження суми простих чисел

# Цикл триває, поки користувач не введе 0
# ( 0 - ознака кінця введення)
while x != 0:
    is_prime = True  # Прапорець: припускаємо,
                     #  що число x є простим
    if x <= 1:
        is_prime = False  # Числа <= 1 не є простими
    else:
        # Натсупний цикл перевіряє наявність дільників 
        # від 2 до  кореня з x включно
        for i in range(2, int(x**0.5) + 1):
            if x % i == 0: # Якщо x ділиться на i без остачі 
                is_prime = False  # то воно не просте, і далі
                break             # можна не перевіряти

    # Якщо число виявилося простим
    if is_prime:
        k += 1  # Збільшуємо кількість простих чисел на 1
        S += x  # Додаємо число до загальної суми

    x = int(input())  # Зчитуємо наступне число для наступної 
                      # ітерації циклу

# Підбиття підсумків після виходу з циклу
if k == 0: # Якщо жодного простого числа не було введено
    print("Простих чисел немає") 
else:
    # Обчислюємо та виводимо середнє арифметичне (сума/кількість)
    print("Середнє арифметичне простих:", S / k)
\end{py}
Програма почергово приймає цілі числа, поки не буде введено 0 (сигнал зупинки). Для кожного введеного числа перевіряється, чи є воно простим (тобто більшим за 1 і ділиться лише на 1 та саму себе).  Якщо число виявилося простим, програма збільшує лічильник $k$ на 1 та додає це число до загальної суми $S$. Після введення 0 програма ділить суму простих чисел на їхню кількість ($S$ / $k$) або повідомляє, що простих чисел не було введено. 
\begin{py}{Приклад 8. Перевірка наявності цифри 4 у натуральному числі}
N = int(input('N = '))  # Зчитуємо натуральне число
temp = N                # Створюємо тимчасову копію числа,
# щоб не змінювати оригінальне N

# Цикл триває, поки у тимчасовій змінній залишаються цифри 
#(число більше 0)
while temp > 0:
    x = temp % 10   # Отримуємо останню цифру числа 
                    #(остача від ділення на 10)

    if x == 4:      # Перевіряємо, чи дорівнює ця цифра 4
        print('Є 4')#Виводимо повідомлення, що цифру 4 знайдено
        break   # Достроково перериваємо цикл, 
                # оскільки перевіряти далі немає потреби
    #Відкидаємо останню цифру 
    temp //= 10) #(цілочисельне ділення на 10)

else: # Блок else виконується лише тоді,
   # коли цикл while завершився природним шляхом
   # (тобто temp стало <= 0), і не був перерваний break
   print('Немає 4')
\end{py}
Користувач вводить число N. Програма створює його копію temp, щоб зберегти оригінальне значення. Отримує останню цифру числа за допомогою temp \% 10.   Якщо цифра дорівнює 4, виводиться 'Є 4', і цикл негайно зупиняється (break).   Якщо цифра не 4, вона відкидається за допомогою цілочисельного ділення temp на 10. Якщо цикл закінчився і цифру 4 так і не знайшли, спрацьовує блок else, який виводить 'Немає 4'.   
\begin{py}{Приклад 9. Алгоритм Евкліда пошуку найбільшого спільного дільника (НСД)}
a = int(input("a = "))
b = int(input("b = "))

while a > 0 and b > 0:
    if a > b:
        a %= b
    else:
        b %= a

print("НСД =", a + b)
\end{py}
	\section*{Контрольні питання  }

\begin{enumerate}
	\item Який синтаксис циклу \texttt{while} у Python?
	\item У чому полягає відмінність між циклами \texttt{for} та \texttt{while}?
	\item Коли доцільно використовувати цикл \texttt{while}, а коли цикл \texttt{for}?
	\item Що означає «цикл із передумовою»?
	\item Що відбувається, якщо умова циклу \texttt{while} від самого початку має значення \texttt{False}?
	\item Як працюють оператори \texttt{break}, \texttt{continue}, \texttt{pass} у циклі \texttt{while}?
	\item Для чого використовується блок \texttt{else} у циклі \texttt{while}?
	\item Що таке «зациклювання»? Які основні причини його виникнення?
	\item Як уникнути нескінченного циклу у програмі?
	\item Наведіть приклад задачі, де цикл \texttt{while} зручніший за цикл \texttt{for}.
\end{enumerate}



\section*{Практичні завдання}

\begin{enumerate}
	\item Надрукувати числа від 0 до 9 за допомогою циклу \texttt{while}.
	\item Обчислити суму чисел від 1 до $n$, використовуючи цикл \texttt{while}.
	\item Запитувати у користувача пароль доти, доки він не введе правильний.
	\item Вивести таблицю квадратів чисел від 1 до 10.
	\item Реалізувати програму, яка знаходить найбільший спільний дільник двох чисел (алгоритм Евкліда).
	\item Написати програму, яка виводить усі цифри введеного числа по одній.
	\item Обчислити факторіал числа $n$ за допомогою циклу \texttt{while}.
	\item Запитувати у користувача додатне число доти, доки він не введе коректне значення.
	\item Обчислити середнє арифметичне простих чисел, що вводяться користувачем, до введення нуля.
	\item Реалізувати програму, яка перевіряє, чи містить введене число цифру нуль.
	\item Обчислити суму введених чисел до введення нуля (стоп‑сигнал).
	\item Написати нескінченний цикл \texttt{while True}, який виводить повідомлення, і перервати його комбінацією клавіш \texttt{Ctrl+C}.
\end{enumerate}

